\documentclass[11pt]{article}
\usepackage[a4paper,margin=30mm]{geometry}
\usepackage{amsmath,amssymb,amsthm}
\usepackage[hidelinks]{hyperref}
\hypersetup{
  pdftitle={Two-sided computer-assisted progress on Erdos Problem 302},
  pdfauthor={Dmitry Khanukov}
}

\newtheorem{theorem}{Theorem}
\newtheorem{proposition}[theorem]{Proposition}
\newtheorem{lemma}[theorem]{Lemma}
\theoremstyle{definition}
\newtheorem{definition}[theorem]{Definition}
\theoremstyle{plain}
\newcommand{\N}{\mathbb N}
\newcommand{\Q}{\mathbb Q}

\title{Two-sided computer-assisted progress on Erd\H{o}s Problem 302}
\author{Dmitry Khanukov}
\date{Version 0.2.0-preprint, 10 September 2026\\[0.5ex]
\textbf{Preliminary and unrefereed}}

\begin{document}
\maketitle

\begin{abstract}
Let $f(N)$ denote the largest size of a subset of $\{1,\ldots,N\}$
containing no distinct $a,b,c$ with $1/a=1/b+1/c$.  We prove two partial
results.  First, using the structured Problem 301 construction of Della
Pietra, we show via an odd-quarter padding lemma that there is an absolute
$\delta>0$ such that
\[
  f(N)\geq(5/8+\delta)N
\]
for all sufficiently large $N$.  The required Problem 301 input comes from a
pinned, kernel-checked but unrefereed external Lean development.  Second, a
hierarchical exact rational
certificate on the nontrivial divisors of $139708800$ gives
\[
  \limsup_{N\to\infty}\frac{f(N)}N
  \leq\frac{140803024}{163562355}
  \approx 0.860852266403232.
\]
The upper finite certificate is checked by a dependency-free exact
verification path, and the upper certificate, semantic bridge, finite
omission argument, and asymptotic endpoint are checked end to end in Lean 4.
The upper bound does not depend on the Della Pietra developments.  The lower
constant is qualitative and non-explicit, and the external developments and
this manuscript are unrefereed. Neither result closes Problem 302.
\end{abstract}

\section{Statement and prior context}

Erd\H{o}s and Graham asked for the asymptotic size of a largest set avoiding
the displayed two-tail reciprocal relation \cite{ErdosGraham1980,Erdos302}.
The elementary interval construction gives the density-$5/8$ lower bound
recorded in the Problem 302 remarks and attributed there to Cambie
\cite{Cambie302}.  Van Doorn obtained
\[
 f(N)\leq(9/10+o(1))N
\]
using disjoint multiplicative dilates \cite{VanDoorn302}.  A separate
five-point argument of van Doorn recorded on the official Problem 301 page
\cite{Erdos301} gives the transferable bound
\[
 f(N)\leq\left(\frac{25}{28}+o(1)\right)N
\]
for Problem 302.  More precisely, the configuration
\[
 \{2t,3t,4t,6t,12t\}
\]
is already constrained by the three two-tail identities
\[
 \frac1{2t}=\frac1{3t}+\frac1{6t},\qquad
 \frac1{3t}=\frac1{4t}+\frac1{12t},\qquad
 \frac1{4t}=\frac1{6t}+\frac1{12t}.
\]
Every one-point deletion of the five-point block leaves at least one of these
relations, so a full block forces two omissions; the truncated block
$\{2t,3t,4t,6t\}$ forces one via the first relation.  Hence the associated
disjoint-dilate count applies verbatim to Problem 302.  This observation
transfers only that finite two-tail configuration argument; it does not treat
the full all-tail statement of Problem 301 as an upper-bound theorem for
Problem 302.

Wang subsequently proved the stronger bound $667/806$ for Problem 301 using
the nontrivial divisors of $720$, exact prefix independence numbers, and a
disjoint family of multiplier dilates \cite{Wang301_2026}.  That theorem does
not directly give an upper bound for Problem 302: avoiding all tail lengths is
a stricter condition than avoiding only two-tail relations, so the Problem
301 extremal function is at most the Problem 302 extremal function.  However,
Wang's finite tile and multiplier family can be specialized to Problem 302 by
discarding every hyperedge except the two-tail triples.  An independent exact
recomputation gives the first prefix thresholds at cover levels $1$ through
$11$ as
\[
 6,12,24,30,40,48,90,120,144,180,240.
\]
Their reciprocal sum is $293/720$, while Wang's multiplier family has density
$120/403$.  The same disjoint-prefix count therefore gives the derived
Problem 302 bound
\[
 f(N)\leq\left(1-\frac{120}{403}\frac{293}{720}+o(1)\right)N
      =\left(\frac{2125}{2418}+o(1)\right)N.
\]
This specialization, rather than the $667/806$ theorem itself, is the
transferable comparison.  The accompanying dependency-free standard-library
verifier regenerates all divisors and two-tail edges, exhaustively computes
the prefix cover numbers, and checks the displayed rational arithmetic.

A publicly available anonymous partial proof posted independently in July 2026
\cite{Partial302Pastebin2026} used the disjoint-dilate architecture and obtained
\[
 f(N)\leq\left(\frac{373}{420}+o(1)\right)N.
\]
The present work further strengthens the finite-block method using
hierarchical exact rational packings and obtains
\[
 \limsup_{N\to\infty}\frac{f(N)}N
 \leq\frac{140803024}{163562355}.
\]
Numerically and exactly, the relevant upper bounds are ordered as
\[
 \frac{140803024}{163562355}<\frac{2125}{2418}
 <\frac{373}{420}<\frac{25}{28}<\frac9{10}.
\]
It does not prove that $f(N)/N$ converges.

Our contribution to the lower bound is an odd-quarter padding lemma converting
Della Pietra's structured Problem 301 witness into a Problem 302 construction
of density strictly greater than $5/8$.  The upper bound is independent of the
Della Pietra developments.

\begin{theorem}\label{thm:main}
There is an absolute $\delta>0$ and an $N_0$ such that, for every $N\geq N_0$,
\[
 f(N)\geq(5/8+\delta)N.
\]
Moreover,
\[
 \limsup_{N\to\infty}\frac{f(N)}N
 \leq\frac{140803024}{163562355}.
\]
\end{theorem}

\section{The lower bound}

The following proposition is imported from the pinned Della Pietra Problem
301 development \cite{DellaPietra301} and is not original to this work.  Put
\[
 P_L=\prod_{p<L}p,
 \qquad
 \rho_L=\frac{\varphi(P_L)}{P_L}
       =\prod_{p<L}\left(1-\frac1p\right).
\]
The product includes the prime $2$.

\begin{proposition}[structured Problem 301 witness]\label{prop:structured}
For some fixed $L$, some fixed auxiliary parameters, and all sufficiently
large $N$, there is a set $C_N\subseteq\{1,\ldots,N\}$ such that:
\begin{enumerate}
\item $C_N$ contains no all-distinct reciprocal relation of any tail length
      at least two;
\item every $n\in C_N$ satisfies either
\[
 N<3n,\qquad 2n<N,\qquad n\ \text{odd},
\]
or $N\leq2n$;
\item
\[
 |C_N|\geq\left(\frac12+\frac{\rho_L}{24}\right)N.
\]
\end{enumerate}
\end{proposition}

All subsequent arguments in this section are elementary except for this
imported structured input.  The new ingredient is
Lemma~\ref{lem:padding}, the odd-quarter padding lemma.

The exact Lean interfaces and immutable dependency revisions used for this
proposition are recorded in the repository.  The underlying analytic work
uses the centered-tail library developed for Problem 327
\cite{DellaPietra327}, whose human manuscript in turn invokes standard
analytic estimates \cite{Tenenbaum2015,BretecheTenenbaum2026}.  These external
developments are unrefereed.

The proposition enters the Lean dependency closure as an exact pinned proof
term, not as a hypothesis of Theorem~\ref{thm:main}.  Mathematical confidence
in the lower result nevertheless depends on the correctness and intended
interpretation of that unrefereed upstream formalization.  The verification
and trust-boundary details are given in the final section.  The upper bound of
Theorem~\ref{thm:main} does not depend on these external developments.

Define the odd quarter
\[
 O_N=\{n\in\{1,\ldots,N\}: n\text{ is odd and }4n\leq N\},
 \qquad A_N=C_N\cup O_N.
\]

\begin{lemma}[odd-quarter padding]\label{lem:padding}
The set $A_N$ contains no distinct $a,b,c$ with
$1/a=1/b+1/c$.
\end{lemma}
\begin{proof}
Assume such a triple exists and relabel the two tails so that $b<c$.
Every reciprocal relation has $a<b<c$.  Moreover $b<c$ gives
$1/a=1/b+1/c>2/c$, so $2a<c\leq N$.

Suppose first that $a\in C_N$.  The top-half alternative is impossible, so
$N<3a$.  Since $b,c>a>N/3>N/4$, neither tail lies in $O_N$.  Thus all three
denominators lie in $C_N$, contradicting its Problem 301 admissibility.

Now suppose that $a\in O_N$.  The head is odd.  From
\[
 bc=a(b+c)
\]
both tails are even: two odd tails make the two sides have opposite parity,
as does one odd and one even tail.  Therefore $b,c\notin O_N$ and both tails
belong to $C_N$.  Their low-band alternative is odd, so both lie in the top
half:
\[
 N\leq2b,\qquad N\leq2c.
\]
Since $b\ne c$, the two tails cannot both equal $N/2$.  Thus at least one of
the reciprocal inequalities $1/b\leq2/N$ and $1/c\leq2/N$ is strict, and
hence
\[
 \frac1b+\frac1c<\frac4N.
\]
On the other hand $4a\leq N$ gives $1/a\geq4/N$, a contradiction.
\end{proof}

The map $j\mapsto2j+1$ embeds $\{0,\ldots,\lfloor N/8\rfloor-1\}$ into
$O_N$, so
\[
 |O_N|\geq\lfloor N/8\rfloor\geq N/8-1.
\]
The two sets are disjoint: every element of $O_N$ is at most $N/4$, whereas
each alternative in Proposition~\ref{prop:structured} lies above $N/3$.
Consequently
\[
 |A_N|\geq
 \left(\frac58+\frac{\rho_L}{24}\right)N-1.
\]
Since $\rho_L>0$, increasing $N_0$ absorbs the additive loss and gives
Theorem~\ref{thm:main}'s lower statement with
\[
 \delta=\frac{\rho_L}{48}>0.
\]
No numerical value of $L$ or $\delta$ is asserted.

The local Lean project also defines $f_{302}(N)$ as the finite maximum of the
cardinalities of triple-free subsets of $\{1,\ldots,N\}$, proves that every
such witness is bounded by this maximum, and derives the lower statement of
Theorem~\ref{thm:main} in that conventional notation.  Separate examples
verify that $\{2,3,6\}$ and every full interval with $N\geq6$ violate the
literal predicate, so the formal statement is not vacuously true.

Readers checking the formalization should go to a single declaration, namely
\[
 \texttt{Erdos302Lower.erdos302\_f302\_lower\_five\_eighths\_plus}
\]
in the file \texttt{lower-lean/Erdos302Lower/Maximum.lean}.  Unfolding its
type
\[
 \texttt{Erdos302MaximumLowerConclusion}
\]
gives
\[
 \exists\,\delta>0\ \exists\,N_0\ \forall N\geq N_0:\quad
 \left(\frac58+\delta\right)N\leq f_{302}(N).
\]
The unfolded statement mentions the single project definition \texttt{f302},
whose own definition uses \texttt{admissibleSubsets} and
\texttt{NoUnitFractionTriple}; those three are reproduced in the repository
README and are the entire vocabulary required.  The proof modules culminate
in this declaration, while the validation and axiom-audit modules are
independent of it.

\section{The hierarchical finite certificate}

Let
\[
 Q=2^7 3^4 5^2 7^2 11=139708800,
 \qquad D=\operatorname{Div}(Q)\setminus\{1\}.
\]
The reciprocal-triple hypergraph $H$ has vertex set $D$ and one edge for each
unordered all-distinct triple satisfying $bc=a(b+c)$.  The verifier
regenerates $719$ vertices and $12675$ edges from these definitions.

Throughout, a \emph{vertex cover} of a hypergraph is a set of vertices
meeting every edge.  A triple-free subset of $D$ is exactly the complement of
a vertex cover of $H$, so lower bounds for cover numbers are what the upper
bound needs.  Computing those numbers directly on $H$ is infeasible, so we
bound them from below by a weighted family of subsets whose individual cover
requirements are already known.  The next definition names that device; the
two families we actually use are constructed in Lemmas~\ref{lem:edge}
and~\ref{lem:scaling}.

For $W\subseteq D$ let $H[W]$ denote the induced sub-hypergraph, whose edges
are the edges of $H$ contained in $W$, and let $\tau(\cdot)$ denote the
minimum size of a vertex cover.

\begin{definition}[configuration, demand, support]\label{def:configuration}
A \emph{configuration} is a pair $(S,r)$ consisting of a subset $S\subseteq
D$, called its \emph{support}, and an integer $r\geq0$, called its
\emph{demand}, such that
\[
 \tau\bigl(H[S]\bigr)\geq r .
\]
\end{definition}

The demand is therefore a proved property of the pair, not a label attached
to it: exhibiting a configuration requires proving the displayed inequality.
The point of the definition is that a demand established once remains valid
inside every larger vertex set, which is what lets a single family serve all
the prefixes used below.

\begin{lemma}[configuration packing]\label{lem:packing}
Let $W\subseteq D$, let $\mathcal G$ be a finite family of configurations
with $S\subseteq W$ for every $(S,r)\in\mathcal G$, and let
$y_{(S,r)}\geq0$ be rational weights satisfying the load condition
\[
 \sum_{(S,r)\in\mathcal G:\,v\in S}y_{(S,r)}\leq1\qquad(v\in W).
\]
Then every vertex cover $C$ of $H[W]$ satisfies
\[
 |C|\geq\sum_{(S,r)\in\mathcal G}r\,y_{(S,r)}.
\]
\end{lemma}
\begin{proof}
Fix $(S,r)\in\mathcal G$.  Every edge of $H[S]$ is an edge of $H[W]$ and is
contained in $S$, so it is met by $C$ inside $S$; hence $C\cap S$ is a vertex
cover of $H[S]$ and $|C\cap S|\geq\tau(H[S])\geq r$.  Exact double counting
now gives
\[
 \sum_{(S,r)}r\,y_{(S,r)}
 \leq\sum_{(S,r)}|C\cap S|\,y_{(S,r)}
 =\sum_{v\in C}\ \sum_{(S,r):\,v\in S}y_{(S,r)}
 \leq|C|.\qedhere
\]
\end{proof}

The first family of configurations is immediate.

\begin{lemma}[edge configurations]\label{lem:edge}
For every edge $e$ of $H$, the pair $(e,1)$ is a configuration.
\end{lemma}
\begin{proof}
The hypergraph $H[e]$ has $e$ among its edges, so a vertex cover of $H[e]$ is
nonempty and $\tau(H[e])\geq1$.
\end{proof}

The second family transports exactly computed cover numbers from a smaller
divisor set.  Put
\[
 Q_0=3360=2^5\cdot3\cdot5\cdot7,
 \qquad D_0=\operatorname{Div}(Q_0)\setminus\{1\},
\]
so that $|D_0|=47$, and let $H_0$ denote the reciprocal-triple hypergraph on
$D_0$, defined by the same relation $bc=a(b+c)$.  For $x\geq1$ write
\[
 D_{0,\leq x}=\{d\in D_0:d\leq x\}.
\]
Because $D_0$ is small, the exact vertex-cover number
$\tau\bigl(H_0[D_{0,\leq x}]\bigr)$ of every such prefix is computed by
exhaustive search; a separate verifier performs this computation
independently of the main certificate.  These exact numbers rise from $0$ to
$21$, and for $1\leq k\leq21$ we define
\[
 s_k=\min\bigl\{t\in D_0:\tau\bigl(H_0[D_{0,\leq t}]\bigr)\geq k\bigr\},
\]
which yields
\[
 \begin{aligned}
 (s_1,\ldots,s_{21})=(
 &6,12,24,30,40,42,56,84,96,105,120,\\
 &140,210,224,240,336,420,560,1120,1680,3360).
 \end{aligned}
\]
We call $D_{0,\leq s_k}$ the \emph{$k$-th certified base prefix}.  Its exact
cover number is $k$, not merely at least $k$.  Indeed, consecutive prefixes
of $D_0$ differ by a single vertex, and adjoining one vertex $v$ raises the
cover number by at most one, since if $C$ covers a prefix $W$ then
$C\cup\{v\}$ covers $W\cup\{v\}$: every edge of $H_0[W\cup\{v\}]$ either lies
in $W$, and is met by $C$, or contains $v$.  By minimality of $s_k$ the
preceding prefix has cover number at most $k-1$, so the prefix
$D_{0,\leq s_k}$ has cover number at most $k$, hence exactly $k$.  The
exhaustive verifier computes these numbers directly and checks the equality.
This replaces an infeasible computation on $H$ by $21$ feasible ones on
$H_0$.

\begin{lemma}[scaled base prefixes]\label{lem:scaling}
Let $m$ divide $Q/Q_0$ and let $1\leq k\leq21$.  Then
\[
 \bigl(m\cdot D_{0,\leq s_k},\ k\bigr)
\]
is a configuration of $H$.
\end{lemma}
\begin{proof}
Write $S=m\cdot D_{0,\leq s_k}$.  If $d\mid Q_0$ and $m\mid Q/Q_0$, then
$md\mid Q$ and $md\neq1$, so $S\subseteq D$, and $d\mapsto md$ is injective.
It also preserves the defining relation in both directions, since
\[
 \frac1a=\frac1b+\frac1c
 \iff
 \frac1{ma}=\frac1{mb}+\frac1{mc},
\]
so it maps every edge of $H_0[D_{0,\leq s_k}]$ to an edge of $H$ contained in
$S$, that is, to an edge of $H[S]$.  Now let $C$ be a vertex cover of $H[S]$
and put
\[
 C_0=\{d\in D_{0,\leq s_k}:md\in C\}.
\]
Given an edge $e_0$ of $H_0[D_{0,\leq s_k}]$, its image $me_0$ is an edge of
$H[S]$, so $C$ contains some $md\in me_0$, whence $d\in C_0\cap e_0$.  Thus
$C_0$ is a vertex cover of $H_0[D_{0,\leq s_k}]$, and injectivity gives
\[
 |C|\geq|C_0|\geq\tau\bigl(H_0[D_{0,\leq s_k}]\bigr)=k,
\]
so $\tau(H[S])\geq k$.
\end{proof}

The verifier does not take the edge correspondence used in this proof on
trust: for every multiplier $m$ and every certified base prefix, it checks
that each mapped triple actually occurs in the independently regenerated edge
set of $H$.

Since
\[
 \frac{Q}{Q_0}=41580=2^2\cdot3^3\cdot5\cdot7\cdot11
\]
has $3\cdot4\cdot2\cdot2\cdot2=96$ divisors, Lemma~\ref{lem:scaling} supplies
$96\cdot21=2016$ configurations, which we call \emph{hierarchical gadgets}.
Together with the $12675$ edge configurations of Lemma~\ref{lem:edge}, they
form the family $\mathcal G$ used below, of size
\[
 12675+2016=14691.
\]

For $x\geq1$ write
\[
 D_{\leq x}=\{d\in D:d\leq x\}.
\]
We apply Lemma~\ref{lem:packing} with $W=D_{\leq x}$, using those
configurations of $\mathcal G$ whose support lies in $D_{\leq x}$.  Since
cover numbers are integers, a packing whose objective is strictly greater
than $k-1$ certifies $\tau(H[D_{\leq x}])\geq k$.
The committed JSON contains exact rational packings at $271$ prefixes,
producing a nondecreasing sequence of possibly repeated thresholds
\[
 t_1\leq t_2\leq\cdots\leq t_{274}
\]
such that the reciprocal-triple hypergraph induced by $D_{\leq t_k}$ has
vertex-cover number at least $k$.  The verifier regenerates every
configuration and checks every rational load and strict objective inequality.
The resulting exact prefix strength is
\[
 S=\sum_{k=1}^{274}\frac1{t_k}
  =\frac{3251333}{4989600}.
\]

\section{Disjoint multiplier blocks and the upper constant}

Write $Q=\prod_{p^e\parallel Q}p^e$.  For a positive integer $R$, define
\[
 \mathcal M_R=
 \left\{
 m\in\N:
 v_p(m)\in\{0,e+1,2(e+1),\ldots,(R-1)(e+1)\}
 \text{ for every }p^e\parallel Q
 \right\}.
\]
There is no restriction on $v_\ell(m)$ for primes $\ell\nmid Q$.  Explicitly,
the valuation moduli at $2,3,5,7,11$ are $8,5,3,3,2$.

If $m,m'\in\mathcal M_R$ and $md=m'd'$ with $d,d'\in D$, then, for each
$p^e\parallel Q$,
\[
 v_p(d)\equiv v_p(md)=v_p(m'd')\equiv v_p(d')\pmod{e+1}.
\]
Both $v_p(d)$ and $v_p(d')$ lie in $\{0,\ldots,e\}$, so they are equal.
Thus $d=d'$ and then $m=m'$.  Hence the blocks $mD$, for
$m\in\mathcal M_R$, are pairwise disjoint.

For fixed $R$, finite inclusion--exclusion gives
\[
 M_R(x):=\#\{m\in\mathcal M_R:m\leq x\}
       =\rho_Rx+O_{Q,R}(1),
\]
where
\[
 \rho_R=
 \prod_{p^e\parallel Q}
 \left(1-\frac1p\right)
 \sum_{j=0}^{R-1}p^{-j(e+1)}.
\]
Indeed, the density of the condition $v_p(m)=j(e+1)$ is
$(1-1/p)p^{-j(e+1)}$.

Let $A\subseteq\{1,\ldots,N\}$ be triple-free.  In every block, the omitted
vertices form a cover of the applicable divisor prefix.  More precisely, for
$m\in\mathcal M_R$ put
\[
 C_m=\{d\in D_{\leq N/m}:md\notin A\}.
\]
Scaling preserves reciprocal triples, so $C_m$ is a vertex cover of the
hypergraph induced by $D_{\leq N/m}$.  Consequently, with
$\max\varnothing=0$,
\[
 |C_m|\geq\max\{k:mt_k\leq N\}
       =\sum_{k=1}^{274}\mathbf 1_{mt_k\leq N}.
\]
Since the blocks are disjoint and every $t_k\geq2$, this gives the finite
inequality
\[
 \begin{aligned}
 N-|A|
 &\geq
 \sum_{\substack{m\in\mathcal M_R\\m\leq N/2}}|C_m| \\
 &\geq
 \sum_{k=1}^{274}
 \#\left\{m\in\mathcal M_R:m\leq\frac{N}{t_k}\right\}.
 \end{aligned}
\]
Applying this to an extremal $A$ and first letting $N\to\infty$ with $R$
fixed yields
\[
 \liminf_{N\to\infty}\frac{N-f(N)}N
 \geq\rho_R\sum_{k=1}^{274}\frac1{t_k}.
\]
Finally,
\[
 \rho_R\mathrel{\uparrow}
 \rho=\prod_{p^e\parallel Q}
       \frac{1-1/p}{1-p^{-(e+1)}}
     =\frac{23520}{110143}.
\]
Letting $R\to\infty$ therefore gives
\[
 \liminf_{N\to\infty}\frac{N-f(N)}N
 \geq \rho S
 =\frac{22759331}{163562355}.
\]
Thus
\[
 \limsup_{N\to\infty}\frac{f(N)}N
 \leq1-\rho S
 =\frac{140803024}{163562355},
\]
which completes Theorem~\ref{thm:main}.

\section{Verification scope and disclosure}

The upper packing was discovered with AI-assisted search and a
floating-point LP solver.  The published claim depends only on the committed
exact rational certificate and its Lean-kernel-checked proof chain.  The
dependency-free standard-library verifier provides an independent executable
cross-check: it is run under normal and optimized Python, regenerates the
complete finite instance, pins its digest and target sequence, and includes
deliberate corruption tests.  The Lean development checks the finite semantic
data, all 271 natural-number packing certificates, the 274-level ledger, the
prefix-omission certificate, and the asymptotic endpoint.  The final theorem
has no certificate hypothesis; its transitive axiom report contains only
\texttt{propext}, \texttt{Classical.choice}, and \texttt{Quot.sound}.  Its
predicates \texttt{NoUnitFractionTriple} and \texttt{IsMaxNoTripleCard} match
the corresponding Erd\H{o}s-302 definitions in
\texttt{google-deepmind/formal-conjectures} at the pinned source revision (up
to the name of one bound variable), while a separate bridge identifies that
interface with the concrete finite maximum used in the proof.

All 80,181 committed project-local Lean modules have source.  For exact main
commit \texttt{015d376e432040a2a71c7a5689c2fd238779ac2c}, the cache-free
full-project workflow compiled every one of those modules from source,
required the exact inventory of 160,362 \texttt{.olean}/\texttt{.ilean}
outputs, replayed the final declarations, and reproduced the axiom allowlist
in GitHub Actions run
\href{https://github.com/khanukov/erdos302/actions/runs/34676374033}{34676374033}.
The present revision changes no project-local Lean source.
The Lean kernel and toolchain, operating system and hardware, and Mathlib are
disclosed trusted boundaries of that source build.

The lower result imports analytic input from pinned Della Pietra Problem
301/327 Lean developments.  At commit
\texttt{4c365e9ded04f04ecd9a6d89a38f97c529194475}, the separate lower project
built from the exact pins and reproduced its committed transitive axiom report
in GitHub Actions run
\href{https://github.com/khanukov/erdos302/actions/runs/31909273465}{31909273465}.
That historical run used the binary-cache-backed build and predates the
current blocking fresh-replay step. On pushes to \texttt{main}, version tags,
and manual runs, the present workflow additionally runs
\texttt{lake env leanchecker {-}{-}fresh Erdos302Lower} over the complete top
module import closure. Pull-request CI still performs the source audit, exact
build, and axiom comparison, but skips this long replay. Tag and manual runs
provide additional replay evidence; automated publication accepts only the
successful push-to-main run for the exact commit. The upstream results are
imported as proof terms, not as hypotheses.
The transitive axiom transcript contains only \texttt{propext},
\texttt{Classical.choice}, and \texttt{Quot.sound}; no project-local open proof
obligation or Problem 301 hypothesis occurs in the theorem statement.  This
records its axiom closure, not independent review or the pedigree of the
software stack: the exact build deliberately uses prerelease Lean
\texttt{v4.33.0-rc1} and the pinned \texttt{teorth/mathlib4} fork containing
the then-unmerged \texttt{Mathlib.NumberTheory.\allowbreak Mertens} module.
The ordinary build downloads precompiled Mathlib \texttt{.olean} files; fresh replay
kernel-checks their stored proof terms but structurally trusts their
serialization. The checker uses Lean's own kernel and is not an independent
checker. The axiom transcript and fresh replay are separate checks.

As of 15 August 2026, the pinned Problem 327 base has an open
\href{https://github.com/donalddellapietra/erdos-327-proof/pull/1}{upstream
exposition correction}.  It changes no Lean file: the exact identity was
already proved as \texttt{mixed\_odd\_factorCount\_eq}, and the correction
changes no theorem, exponent, parameter choice, or certified margin.  The
lower build therefore remains on the immutable base pin rather than silently
following an open pull request.  A root-project build alone does not establish
the lower statement.
At the pinned revisions, the external repositories have no license file, so
their source is not vendored here.

AI systems assisted with search, code generation, proof audits, and
formalization.  This assistance does not replace author or referee review.
Both the external work and this manuscript are unrefereed.  A fully verified
commit may be merged into the repository as unrefereed engineering work, with
referee corrections applied in follow-up pull requests.  A tagged preliminary
GitHub/Zenodo release and an arXiv preprint may be published before independent
human review to establish a public timestamp, provided that the exact released
commit passes every proof-boundary check and the release is explicitly labelled
preliminary and unrefereed.  Such a release must not claim independent
verification, peer review, or a solution of Problem 302.  Named independent
human mathematical review is required before describing the result as
independently verified.

Original manuscript material is licensed under
\href{https://creativecommons.org/licenses/by/4.0/}{CC BY 4.0}.  Attribution
and third-party exclusions are stated in the accompanying license notice.

\bibliographystyle{plain}
\bibliography{references}

\end{document}
